\documentclass[12pt]{article}

% Layout.
\usepackage[top=1in, bottom=0.75in, left=0.75in, right=0.75in, headheight=1in, headsep=6pt]{geometry}

% Fonts.
\usepackage{mathptmx}

\usepackage[scaled=0.86]{helvet}
\renewcommand{\emph}[1]{\textsf{\textbf{#1}}}

% Misc packages.
\usepackage{amsmath,amssymb,latexsym,amsthm}
\usepackage{graphicx,hyperref}
\usepackage{array}
\usepackage{xcolor}
\usepackage{multicol}
\usepackage{tabularx,colortbl}
\usepackage{enumitem}
\usepackage{soul,multicol}
\usepackage{setspace}


\hypersetup{
    colorlinks=true,
    linkcolor=blue,
    filecolor=magenta,      
    urlcolor=blue,
    pdftitle={Proofs Worksheet},
    pdfpagemode=FullScreen,
    }

\def\mailto#1{\href{mailto:#1}{#1}}

%% special math symbols
\newcommand{\N}{\mathbb{N}}
\newcommand{\Z}{\mathbb{Z}}
\newcommand{\R}{\mathbb{R}}
\newcommand{\Q}{\mathbb{Q}}
\newcommand{\sse}{\subseteq}
\newcommand{\mcp}{\mathcal{P}}
\newcommand{\ol}[1]{\overline{#1}}

%other specials
\newcommand{\be}{\begin{enumerate}}
\newcommand{\ee}{\end{enumerate}}
\newcommand{\se}{\subseteq}
\newcommand{\spe}{\supseteq}
\newcommand{\ds}{\displaystyle}
\newcommand{\lcm}{\text{lcm}}
%\newcommand{\gcd}{\text{gcd}}

% Paragraph spacing
\parindent 0pt
\parskip 6pt plus 1pt
\def\tableindent{\hskip 0.5 in}
\def\ts{\hskip 1.5 em}

%header
\usepackage{fancyhdr}
\pagestyle{fancy} 
\lhead{\large\sf\textbf{MATH 265: Introduction to Mathematical Proofs}}
\rhead{\large\sf\textbf{Worksheet 23: Ch 12 \S 3-4}}

\newcommand{\localhead}[1]{\par\smallskip\textbf{#1}\nobreak\\}%
\def\heading#1{\localhead{\large\emph{#1}}}
\def\subheading#1{\localhead{\emph{#1}}}

\newenvironment{clist}%
{\bgroup\parskip 0pt\begin{list}{$\bullet$}{\partopsep 4pt\topsep 0pt\itemsep -2pt}}%
{\end{list}\egroup}%

\begin{document}
\begin{enumerate}
\item The Pigeonhole Principle
	\begin{enumerate}
	\item (version 1)
	\vspace{0.5in}
	\item (version 2)
	\vspace{0.5in}
	\end{enumerate}
\item  Given any 10-element subset of $\{1,2,3, \cdots, 100\}$ there exist at least two distinct subsets with the same sum.
\vfill
\item Let $B=\large\{ \{0\}, \{1,8\},\{2,7\}, \{3,6\},\{4,5\}\large\}$ be a partition  of $\{0,1,2,3,\cdots,9\}.$ \\ Suppose that $f(n): \N \to B$ maps a natural number to the subset of $B$ that contains its remainder upon division by $9.$
	\begin{enumerate}
	\item  Determine $f(13)$ and $f(28).$
	\vspace{1in}
	\item Prove that if $f(n)=f(m)$, then either $9 \mid (n+m)$ or $9 \mid (n-m)_.$
	\vfill
	\newpage
	\item Fill in the blank below with the \emph{smallest} possible value for which it is true and then prove the statement using PHP v2. \\
	
	{Every} subset of $\N$ with at least \underline{\hspace{.5in}} elements will contain a pair of numbers whose sum or difference is divisible by 9. 
	
	\vfill
	\end{enumerate}
\item Let $A=\{a,b,c\}, B=\{1,2,3\},$ and $C=\{4,5,6,7\}.$ \\
Suppose $f:A \to B,$ $g:B \to C$ and $h: C \to B$ are functions defined below.	\\
$$f=\{(a,1),(b,2),(c,2)\}, \quad g=\{(1,4),(2,5),(3,6)\}, \quad  h=\{(4,1),(5,2),(6,2),(7,3)\}.$$
Describe each expression below by using "potato" pictures, stating explicitly the associated ordered pairs, and establishing the domain, codomain and range.  Note ``$\circ$" represents composition of functions just like in precalculus:
$$g\circ f = g(f(x))$$
	\begin{enumerate}
	\item $g \circ f$
	\vfill
	\item $(h \circ g) \circ f$
	\vfill
	\end{enumerate}
\end{enumerate}
\end{document}










